\begin{table}[!htbp]
\centering
\caption{Register of calculation series and consumed propellant. The final recommendation belongs to v4-F.}
\label{tab:versions}
\small
\setlength{\tabcolsep}{4pt}
\begin{tabularx}{\linewidth}{@{}lXrr@{}}
\toprule
Series & Calculation and role & At 110 t & At 120 t\\
\midrule
v1-B & Baseline load scan, 75 km entry ignition & 107.2 t & ---\\
v1-D2 & 2D reference for the crosswind diagnostic & --- & 113.0 t\\
v1-D3 & 3D crosswind diagnostic & --- & 111.9 t\\
v1-J & Jet-interaction sensitivity cases & --- & 111.9--112.9 t\\
v1-P & Pre-search comparator in the guidance comparison & --- & 111.8 t\\
v3-P & Parameterized guidance, entry ignition near 55 km & 103.3 t & 108.5 t\\
v3-I & Ideal reference in the actuator/area ablation & 109.2 t & 114.8 t\\
\textbf{v4-F} & \textbf{Combined final model; recommended design} & \textbf{106.6 t} & \textbf{112.7 t}\\
\bottomrule
\end{tabularx}
\par\smallskip{\footnotesize Column headings indicate initial total propellant, while entries indicate consumption. A dash denotes no value tabulated for that series. The v1-B entry at 110 t is the displayed phase sum $42.7+64.5=107.2$ t; v1-D3 and v1-J also use displayed phase sums. Suffixes identify separate calculation series, not inferred equivalence of controllers or atmosphere settings. v3-A, v4-R, and v4-L denote the area-only, response/staging-only, and 3 s lag variants within the ablation comparison.}
\end{table}
